% ID: SEL-D0033-OB-E02
% ESTADO: cotejada. Segunda pasada visual y borrador compilado revisados.
% Fuente: Selectividad/matematicas_aplicadas/titular_septiembre/Examen-Matemáticas Aplicadas a las Ciencias Sociales.pdf
% SHA-256: 0f95b6aefa78f1a13e2115b281338bd807177671c303c380180957fe03f0a605
% Página: 2; opción B; ejercicio original 2.

El beneficio, en miles de euros, que ha obtenido una almazara a lo largo de 50 años viene dado por la expresión
\[
B(t)=\begin{cases}
-0.04t^2+2.4t & 0\leq t<40\\[4pt]
\dfrac{40t-320}{t} & 40\leq t\leq50
\end{cases}
\]
donde $t$ es el tiempo transcurrido.
\begin{enumerate}[label=\alph*)]
\item \textbf{(1 punto)} Estudie la continuidad y la derivabilidad de la función $B(t)$ en el intervalo $[0,50]$.
\item \textbf{(1 punto)} Estudie la monotonía de la función $B(t)$ y determine en qué momento fueron mayores los beneficios de la almazara, así como el beneficio máximo.
\item \textbf{(0.5 puntos)} Represente la gráfica de la función y explique la evolución del beneficio.
\end{enumerate}
